% Figure: Cassegrain two-mirror telescope. A concave paraboloidal primary of
% focal length f1 gathers a collimated bundle and directs it toward its focus;
% a convex hyperboloidal secondary intercepts the converging cone short of that
% focus and relays it back through a central hole in the primary to the final
% Cassegrain focus, lengthening the effective focal length within a short tube.
\begin{figure}[htbp]
\centering
\begin{tikzpicture}[>=Stealth, scale=1.0]
% optical axis
\draw[enggray, dashed] (-0.6,0) -- (7.2,0);
\node[enggray, font=\scriptsize, anchor=west] at (7.2,0) {axis};
% primary mirror: concave arc opening to the left, centred on axis at x=6
\draw[very thick, engblue] (6,-2.0) .. controls (5.2,-0.7) and (5.2,0.7) .. (6,2.0);
\node[engblue, font=\scriptsize, anchor=west] at (6.05,-1.7) {primary $f_1$};
% central hole in primary
\draw[engblue, thick] (5.72,0.35) -- (6.0,0.35);
\draw[engblue, thick] (5.72,-0.35) -- (6.0,-0.35);
% secondary mirror: small convex arc near x=2, opening to the right
\draw[very thick, engred] (2.0,-0.55) .. controls (2.35,-0.2) and (2.35,0.2) .. (2.0,0.55);
\node[engred, font=\scriptsize, anchor=south] at (2.0,0.6) {secondary};
% incoming collimated rays (two, symmetric)
\draw[engorange, ->, thick] (-0.6,1.5) -- (5.55,1.5);
\draw[engorange, ->, thick] (-0.6,-1.5) -- (5.55,-1.5);
% reflected off primary toward common focus behind secondary (near x=2.6 on axis)
\draw[engorange, thick] (5.55,1.5) -- (2.28,0.42);
\draw[engorange, thick] (5.55,-1.5) -- (2.28,-0.42);
% intercepted by secondary and relayed through hole to Cassegrain focus at x=6.9
\draw[engamber, thick] (2.28,0.42) -- (5.86,0.05);
\draw[engamber, thick] (2.28,-0.42) -- (5.86,-0.05);
\draw[engamber, ->, thick] (5.86,0.05) -- (6.95,0.0);
\draw[engamber, ->, thick] (5.86,-0.05) -- (6.95,0.0);
% Cassegrain focus
\filldraw[engblack] (6.95,0) circle [radius=1.4pt];
\node[font=\scriptsize, anchor=north west] at (6.95,-0.05) {Cassegrain focus};
% back-focal-distance bracket
\draw[<->, enggray] (6.0,-2.15) -- (6.95,-2.15);
\node[enggray, font=\scriptsize, anchor=north] at (6.5,-2.2) {BFD};
% mirror separation bracket
\draw[<->, enggray] (2.0,2.0) -- (6.0,2.0);
\node[enggray, font=\scriptsize, anchor=south] at (4.0,2.0) {$d$};
\end{tikzpicture}
\caption[Cassegrain two-mirror telescope]{A Cassegrain telescope folds a long
effective focal length into a short tube. The concave paraboloidal primary of
focal length $f_1$ converges the incoming collimated bundle toward its own focus;
a convex hyperboloidal secondary intercepts the cone short of that focus and
relays it back through a central hole in the primary to the Cassegrain focus. The
secondary magnifies the primary's cone by $m = f_2'/f_2$, so the system focal
length is $f_{\mathrm{sys}} = m f_1$, several times the primary's, while the
physical tube length stays near the primary focal length. The two conic constants
are fixed by the requirement that the pair be free of spherical aberration, which
is why a classical Cassegrain pairs a paraboloid with a hyperboloid.}
\label{fig:vol04ch11:cassegrain-geometry}
\end{figure}
